\section{Methodology}
\label{sec:methodology}

This section operationalizes the research questions defined in
Section~I into a concrete experimental protocol. The methodology is
structured as \textbf{sequential blocks of activities}, where each
block depends on the completion of the preceding one. Following this
set of activities enables the systematic collection of the elements
required to achieve the research objectives. This block-based
structure also ensures that any independent researcher can replicate
the experiment by following the same sequence of steps.

The study integrates two existing evaluation frameworks---IaC-Eval~[1]
for functional correctness and SecLLM~[2] for security smell
detection---into a single automated pipeline that measures both
dimensions over the same generation tasks, under both baseline and
security-oriented prompting conditions.

% ======================================================================
\subsection{Block 1 --- Conceptual Framework}
\label{sec:block1}

Prior work has evaluated functional correctness~[1] and security
quality~[2],~[5] in isolation, on different datasets, with different
models, and under different prompting conditions. While recent studies
have begun to measure both dimensions jointly---most notably
CWEval~[21], which documents a~30-percentage-point gap between
functional correctness and joint functional-and-secure correctness in
general-purpose languages, and GenSIaC~[15], which demonstrates that
security-aware fine-tuning improves both dimensions simultaneously---
these evaluations serve to validate specific proposed generation
methods rather than to characterize the correlation between the two
dimensions across the contemporary model landscape in the IaC domain.
No study has yet jointly characterized functional correctness and
security quality across multiple LLMs over a dedicated IaC benchmark,
treating their correlation---rather than the validation of any single
generation method---as its central object of inquiry.

This study adopts a \textbf{dual evaluation model} that measures
three dimensions simultaneously:

\begin{enumerate}[label=\textbf{Dim \arabic*:}, leftmargin=*]
    \item \textbf{Functional Correctness (\textit{Fc}).} Does the
          generated Terraform code compile and satisfy the semantic
          intent specification encoded in the Rego policy?
    \item \textbf{Security Quality (\textit{Sp}, \textit{Sc}).}
          \textit{Sp} is a binary indicator ($Sp \in \{0,1\}$, where
          $Sp = 1$ denotes a script free of detected smells and $Sp = 0$
          denotes at least one detected smell). \textit{Sc} (Smell Count)
          is the total number of security smells detected per script,
          a non-negative discrete variable.
    \item \textbf{Fc-Sp Intersection.} What is the relationship
          between both dimensions? Does functionally correct code
          tend to be secure, or does a systematic gap exist?
\end{enumerate}

The three research questions are derived directly from this
framework: RQ1 measures the association between Dimension~1 and the
continuous smell count (\textit{Sc}) from Dimension~2; RQ2 measures
the effect of security-oriented prompting on both dimensions; and
RQ3 characterizes the smell types that persist in code that passes
functional validation (i.e., scripts with $Fc = 1$).

% ======================================================================
\subsection{Block 2 --- Experimental Design}
\label{sec:block2}

The study employs a \textbf{two-factor design} with two independent
variables crossed over the same stratified scenario sample: LLM Model
(M, four levels) and Prompting Strategy (P, two levels: P0 and P1).
Both conditions are executed and evaluated.

\subsubsection{Independent Variables}

\textbf{1. LLM Model (M):} Four models are evaluated, all served
locally through Ollama. The selection spans four architectural
families and three specialization profiles (Table~\ref{tab:models}).

\begin{table}[htbp]
\centering
\caption{Evaluated LLM Models}
\label{tab:models}
\begin{tabular}{@{}llll@{}}
\toprule
\textbf{Model} & \textbf{Family} & \textbf{Params} & \textbf{Profile} \\
\midrule
qwen2.5-coder:7b  & Alibaba    & 7B  & Code-specialized \\
llama3.1:8b        & Meta       & 8B  & General-purpose \\
codellama:7b       & Meta       & 7B  & Code-specialized \\
codegemma:7b       & Google     & 7B  & Code-specialized \\
\bottomrule
\end{tabular}
\end{table}

\textbf{2. Prompting Strategy (P):} Two levels are defined and both
are executed:
\begin{itemize}
    \item \textbf{P0 (Baseline):} A system prompt instructing the
          model to act as an expert Terraform engineer, generate
          valid HCL code, and output only the Terraform code. No
          security guidance is provided.
    \item \textbf{P1 (Security-oriented):} An augmented prompt that
          additionally directs the model to: avoid hardcoded
          credentials; apply least-privilege IAM policies; restrict
          network exposure; enforce encryption in transit; use strong
          cryptography; add integrity checks for downloaded resources;
          and omit placeholder security comments. The instructions
          are derived directly from the eight security smell
          categories operationalized in Block~8.
\end{itemize}

\subsubsection{Dependent Variables}

Three dependent variables capture the experimental outcomes:

\begin{enumerate}
    \item \textbf{Functional Correctness (\textit{Fc}):} A binary
          metric at the script level ($Fc = 1$ if the script passes
          both \texttt{terraform plan} and OPA/Rego semantic intent
          validation; $Fc = 0$ otherwise). At the model level,
          aggregated as \textit{pass@1}, consistent with
          IaC-Eval~[1].
    \item \textbf{Security Smell Prevalence (\textit{Sp}):} A binary
          variable ($Sp = 1$ if no smell is detected in the script;
          $Sp = 0$ if at least one smell is detected).
    \item \textbf{Smell Count (\textit{Sc}):} A non-negative discrete
          count of the total number of security smells detected per
          script. Used as the continuous variable in RQ1 correlations.
    \item \textbf{Smell Taxonomic Type (\textit{St}):} A nominal
          classification according to the eight Terraform-adapted
          rules operationalized from the taxonomy of
          De~Vito et al.~[2] (see Block~8 and
          Table~\ref{tab:smells}).
\end{enumerate}

% ======================================================================
\subsection{Block 3 --- Dataset Description and Preparation}
\label{sec:block3}

The experiment uses the IaC-Eval benchmark~[1], the first dedicated
benchmark for evaluating IaC code generation. It comprises 458
human-curated scenarios covering AWS services across six difficulty
levels, spanning 251 unique AWS resource types. Each scenario provides
a natural language description of the desired infrastructure, a Rego
policy specification for semantic intent validation, and a reference
Terraform output.

\textbf{Role of AWS in this study.} AWS is the \textit{target cloud
provider} described in the benchmark scenarios---not a service that
is actively used during the experiment. The generated Terraform code
is designed to provision AWS resources (S3 buckets, IAM roles, VPCs,
etc.), but the evaluation is performed \textbf{entirely offline}.
Dummy AWS credentials and a provider override disable all real API
calls, allowing \texttt{terraform plan} to validate code structure
without contacting AWS services.

\subsubsection{Dataset Structure}

The dataset is stored in CSV format with six columns:
\texttt{Resource}, \texttt{Prompt}, \texttt{Rego intent},
\texttt{Difficulty} (integer 1--6), \texttt{Reference output}, and
\texttt{Intent}.

\subsubsection{Data Cleaning and Preprocessing}

As IaC-Eval is a human-curated benchmark, no extensive data cleaning
was required. The following preprocessing steps were applied:

\begin{enumerate}[leftmargin=*]
    \item \textbf{Null and empty field validation:} No missing values
          were found in the \texttt{Prompt} or \texttt{Rego intent}
          columns.
    \item \textbf{Sequential indexing:} Each scenario was assigned a
          zero-padded identifier (\texttt{question\_0000},
          \texttt{question\_0001}, \ldots, \texttt{question\_0457}).
    \item \textbf{Rego policy compatibility:} Policies written in
          Rego~v0 syntax (without \texttt{import rego.v1}) require
          the \texttt{-{}-v0-compatible} flag in OPA~v1.17. This
          check was automated in the evaluation script.
    \item \textbf{Stratified random subset selection:} A
          \textbf{stratified random sample} of $N = 40$ scenarios
          was drawn proportionally across the six difficulty
          levels (seed = 42), ensuring representation of all
          complexity strata. The resulting sample indices are
          persisted in \texttt{experiment\_sample.json} so that
          \textbf{exactly the same 40 scenarios} are evaluated for
          all models and both prompting conditions~(P0 and P1),
          guaranteeing comparability. Table~\ref{tab:sample_dist}
          reports the sample distribution by difficulty level.
\end{enumerate}

\begin{table}[htbp]
\centering
\caption{Stratified Sample Distribution by Difficulty Level ($N=40$)}
\label{tab:sample_dist}
\begin{tabular}{@{}lrr@{}}
\toprule
\textbf{Difficulty} & \textbf{Population} & \textbf{Sample} \\
\midrule
1 & 45  & 4  \\
2 & 95  & 8  \\
3 & 113 & 10 \\
4 & 58  & 5  \\
5 & 72  & 6  \\
6 & 75  & 7  \\
\midrule
\textbf{Total} & \textbf{458} & \textbf{40} \\
\bottomrule
\end{tabular}
\end{table}

% ======================================================================
\subsection{Block 4 --- Environment Setup}
\label{sec:block4}

Before executing the pipeline, the following tools must be installed
and configured:

\begin{enumerate}[leftmargin=*]
    \item \textbf{Python 3.10+} with dependencies: \texttt{pandas},
          \texttt{requests}, \texttt{pyyaml}, \texttt{scipy}.
    \item \textbf{Ollama} for local LLM inference, with all models
          downloaded via \texttt{ollama pull}.
    \item \textbf{Terraform CLI} for code compilation and plan
          generation.
    \item \textbf{OPA v1.17} for Rego policy evaluation.
    \item \textbf{AWS credentials offline:} Dummy credentials are
          set via environment variables, and a provider override file
          disables all real AWS API calls.
\begin{verbatim}
provider "aws" {
  access_key = "test"
  secret_key = "test"
  skip_credentials_validation = true
  skip_requesting_account_id  = true
  skip_metadata_api_check     = true
}
\end{verbatim}
    \item \textbf{Plugin cache:} \texttt{TF\_PLUGIN\_CACHE\_DIR} is
          configured to avoid re-downloading the AWS provider.
\end{enumerate}

All infrastructure runs locally on an NVIDIA RTX~4070 Laptop GPU
with 8\,GB of VRAM. All four evaluated models (7--8B parameters)
fit in GPU memory and complete a full generation--evaluation--audit
cycle in approximately 1--2 minutes per scenario.

% ======================================================================
\subsection{Block 5 --- Script Construction: Pipeline Orchestration}
\label{sec:block5}

The complete experiment is orchestrated by an automated Python
script (\texttt{pipeline\_v2.py}) that executes three sequential
phases for each model and prompting condition. The script is
publicly available in the replication
package\footnote{\url{https://github.com/edarmijo/secllm-terraform-integration}}.

The script accepts the following configurable parameters:
\begin{itemize}[leftmargin=*]
    \item \texttt{PIPELINE\_MODELS}: comma-separated list of model
          names (must match \texttt{ollama list}).
    \item \texttt{PIPELINE\_N}: number of scenarios per model per
          condition (default: 40).
    \item \texttt{COMMON\_AUDITOR}: the model designated as the
          external security auditor (default:
          \texttt{devstral-small-2:latest}).
\end{itemize}

Key properties of the script:
\begin{itemize}[leftmargin=*]
    \item \textbf{Resumable:} Re-execution resumes from the last
          completed observation.
    \item \textbf{Incremental:} Results are persisted in JSON after
          each observation.
    \item \textbf{Isolated per model and condition:} Each
          model--condition pair has its own output directory
          (\texttt{outputs\_v2/<model>/<condition>/}).
\end{itemize}

% ======================================================================
\subsection{Block 6 --- Phase 1: Code Generation}
\label{sec:block6}

For each scenario and prompting condition, the model receives the
natural language prompt with a system instruction and generates a
Terraform HCL script. The micro-steps are:

\begin{enumerate}[leftmargin=*]
    \item Load the system prompt (P0 or P1, as defined in Block~2).
    \item Concatenate the system prompt with the scenario natural
          language prompt as the user message.
    \item Send the request to the local Ollama API
          (\texttt{http://localhost:11434/api/generate}) with
          temperature = 0.2 and \texttt{keep\_alive} = 10 minutes.
    \item Extract the HCL block (delimited by \texttt{hcl ...}).
    \item Save the extracted code as
          \texttt{outputs\_v2/<model>/<condition>/terraform/%
          question\_XXXX.tf}.
    \item If extraction fails (empty response or timeout of 1800\,s),
          skip and flag for retry.
    \item Purge conversation context before the next task to ensure
          statistical independence.
\end{enumerate}

% ======================================================================
\subsection{Block 7 --- Phase 2: Functional Evaluation}
\label{sec:block7}

Each generated \texttt{.tf} file is evaluated for functional
correctness through two sequential sub-phases:

\subsubsection{Technical Validation}

\begin{enumerate}[leftmargin=*]
    \item Execute \texttt{terraform init -backend=false} in an
          isolated temporary directory.
    \item Execute \texttt{terraform plan -out plan.out -no-color}.
    \item If the plan fails (return code $\neq 0$): classify as
          $Fc = 0$ and record the error message.
    \item If the plan succeeds: serialize the plan to JSON via
          \texttt{terraform show -json plan.out}.
\end{enumerate}

\subsubsection{Intent Validation}

\begin{enumerate}[leftmargin=*]
    \item Write the corresponding Rego policy to a temporary file.
    \item Execute \texttt{opa eval -i plan.json -d policy.rego data}.
    \item If all leaf values in the OPA output are \texttt{true}:
          classify as $Fc = 1$ (functionally correct).
    \item If any leaf value is \texttt{false}: classify as $Fc = 0$.
\end{enumerate}

Results are recorded incrementally in
\texttt{iac\_eval\_results.json}.

% ======================================================================
\subsection{Block 8 --- Phase 3: Security Auditing}
\label{sec:block8}

Security quality is measured using SecLLM~[2], adapted for
Terraform/HCL and local inference.

\subsubsection{Common External Auditor}

To eliminate the confound introduced by self-auditing---whereby the
same model that generates the code also audits it, entangling
detection capability with willingness to report its own
flaws---this study designates a single \textbf{common external
auditor model} (\texttt{devstral-small-2:latest}) that evaluates
the outputs of all four generation models across both prompting
conditions. This design ensures that security smell detection is
consistent and independent of the generating model.

The auditing micro-steps are:

\begin{enumerate}[leftmargin=*]
    \item Load the SecLLM configuration with the eight
          Terraform-specific security rules
          (\texttt{config\_terraform.yaml}).
    \item For each of the eight rules, send the numbered script to
          the common auditor with the specialized COSTAR prompt.
    \item Parse the structured XML response to determine if the
          smell was detected.
    \item Record the detection result (smell name, line number,
          confidence).
    \item Aggregate results per file: binary prevalence (\textit{Sp}),
          total smell count (\textit{Sc}), and taxonomic
          classification (\textit{St}).
\end{enumerate}

Auditing conditions: temperature = 0.0 (deterministic), timeout =
3600\,s per file.

\subsubsection{Security Smell Taxonomy}

The eight rules operationalize eight of the nine categories from the
taxonomy of De~Vito et al.~[2]. The ninth category, Missing Default
Case, is excluded because the HCL language does not use
switch\,/\,case control structures; this smell is specific to
procedural IaC tools such as Chef and Puppet. The eight implemented
rules use the COSTAR prompt framework with few-shot examples and
structured XML output (Table~\ref{tab:smells}).

\begin{table}[htbp]
\centering
\caption{Security Smell Rules and CWE Mapping (8 Terraform Rules)}
\label{tab:smells}
\begin{tabular}{@{}lll@{}}
\toprule
\textbf{Dimension} & \textbf{Smell} & \textbf{CWE} \\
\midrule
\multirow{3}{*}{Credentials}
  & Admin by Default    & CWE-250 \\
  & Empty Password      & CWE-258 \\
  & Hard-coded Secret   & CWE-798 \\
\midrule
\multirow{2}{*}{Network}
  & Unrestricted IP     & CWE-284 \\
  & HTTP without TLS    & CWE-319 \\
\midrule
\multirow{3}{*}{Code Quality}
  & No Integrity Check  & CWE-353 \\
  & Weak Cryptography   & CWE-327 \\
  & Suspicious Comment  & CWE-546 \\
\bottomrule
\end{tabular}
\end{table}

% ======================================================================
\subsection{Block 9 --- SecLLM Adaptation}
\label{sec:block9}

SecLLM originally supported only Ansible, Puppet, and Chef with
cloud-based LLM APIs. Six modifications were required to adapt it
for this study:

\begin{enumerate}[leftmargin=*]
    \item UTF-8 encoding fixes for configuration and script file
          reading on Windows.
    \item Increased API timeout from 120\,s to 600\,s.
    \item Expansion of the retry loop from 1 to 3 attempts.
    \item Handling of \texttt{confidence = None} for local backends
          that do not return logprobs.
    \item Making \texttt{torch} and \texttt{anthropic} imports
          optional.
    \item Authoring eight Terraform-specific security smell detection
          rules using the COSTAR prompt framework.
\end{enumerate}

All modifications are documented in the replication package.

% ======================================================================
\subsection{Block 10 --- Statistical Analysis Plan}
\label{sec:block10}

Given that the primary metrics are zero-bounded and likely
asymmetric, the analysis uses non-parametric methods. Results are
reported in Section~IV.

\subsubsection{RQ1 --- Correlation}

Spearman's $\rho$ and Kendall's $\tau$ quantify the monotonic
association between $Fc$ (binary) and $Sc$ (smell count, continuous
discrete) across all scripts in the P0 condition. Spearman's $\rho$
also serves as the effect size measure for this test.

\subsubsection{RQ2 --- Prompting Effect}

The Wilcoxon signed-rank test (paired, two-sided) compares $Sc$
and $Fc$ between P0 and P1 on the same $N = 40$ scenarios for each
model. Effect size is quantified by the rank-biserial correlation
$r = W / [n(n+1)/2]$, where $W$ is the Wilcoxon statistic and $n$
is the number of non-zero differences. Multiple comparison correction
uses \textbf{Holm-Bonferroni} (family: RQ1 and RQ2 confirmatory
hypotheses).

\subsubsection{Model Comparison (Exploratory)}

The Kruskal-Wallis test with $\eta^2 = (H - k + 1)/(n - k)$ as
effect size is applied to compare $Sc$ distributions across the
four models. Pairwise Mann-Whitney~$U$ tests with rank-biserial
$r$ as effect size are used for post-hoc comparisons. Multiple
comparison correction for this exploratory family uses
\textbf{Benjamini-Hochberg} FDR control, which is appropriate for
discovery-oriented analyses.

\subsubsection{RQ3 --- Smell Type Distribution}

The analysis is restricted to scripts with $Fc = 1$. Given the
typically small number of such scripts, the analysis is
\textbf{descriptive}: frequency counts and proportions for each
smell type are reported. Fisher's Exact Test is applied only when
cell counts are sufficient ($\geq 5$) and corrected with
Benjamini-Hochberg FDR (exploratory family).

\subsubsection{Multiple Correction Summary}

Two separate families of hypotheses are defined:
\begin{itemize}
    \item \textbf{Confirmatory family} (RQ1 and RQ2): corrected with
          Holm-Bonferroni sequential correction (FWER control).
    \item \textbf{Exploratory family} (model comparisons and smell
          type comparisons in RQ3): corrected with
          Benjamini-Hochberg FDR control.
\end{itemize}

% ======================================================================
\subsection{Block 11 --- Methodology Diagram}
\label{sec:block11}

Figure~\ref{fig:diagrama} presents the complete methodology as a
block diagram, showing the sequential flow of activities from
stratified dataset sampling through parallel P0/P1 code generation,
functional evaluation, common-auditor security auditing, and
statistical analysis.

\begin{figure*}[htbp]

\end{figure*}

% ======================================================================
\subsection{Block 12 --- Threats to Validity}
\label{sec:block12}

\textbf{Construct Validity.} Surface-level metrics such as BLEU or
CodeBERTScore lack validity in the IaC domain, where syntactically
different configurations can be functionally equivalent~[1]. This
study uses \texttt{terraform plan} and OPA/Rego as ground-truth
oracles for functional correctness, and the SecLLM semantic analyzer
for security smell detection. The binary nature of $Fc$ (pass/fail)
may obscure partial correctness; however, this design follows the
IaC-Eval~[1] protocol and enables direct comparison with prior work.

\textbf{Internal Validity.} The use of a single common external
auditor (\texttt{devstral-small-2:latest}) eliminates the
self-auditing confound present in prior versions of this experiment,
where the generating model was also the auditor. The remaining
internal validity concern is that the auditor model may itself have
systematic detection biases; this is mitigated by selecting a
code-specialized model with demonstrated security smell detection
capability. Temperature = 0.0 (deterministic) ensures reproducibility
of audit decisions. Statistical independence between scenarios is
enforced by purging conversation context between tasks.

\textbf{External Validity.} The evaluation is limited to
Terraform/HCL and AWS services. Results may not generalize to other
IaC tools (Ansible, Pulumi, CloudFormation) or cloud providers.
The use of open-weight 7--8B models through Ollama ensures
reproducibility and enables community replication on consumer
hardware, but results may differ for larger proprietary frontier
models (e.g., GPT-4o, Claude Opus).

\textbf{Conclusion Validity.} With $N = 40$ scenarios per model
per condition and four models, the pooled sample for RQ1 is
$N = 160$ observations, providing adequate statistical power for
non-parametric tests. The Wilcoxon signed-rank test for RQ2 operates
on $N = 40$ paired observations per model, which is above the
minimum recommended for reliable inference. The stratified random
sample ensures representation of all six difficulty levels
(Table~\ref{tab:sample_dist}).
